\section{The strategy text, and the hosts checked}
\label{app:strategy}

This appendix holds the part of Study~3 that concerns the strategy level of the skill, a text that teaches the agent
how to organise its tool calls, and the checks of the proxy in MCP hosts (Section~\ref{sec:s3pointer} summarises
them). The hypotheses are numbered as in Table~\ref{tab:s3hyp}, where every outcome is given.

\subsection{E1: the strategy text on the first model}
\label{sec:s3e1strategy}

\paragraph{The strategy alone changed nothing measurable on this model.} Against the arm without it, the median of S
stayed inside R's range for turns, tool calls, billed input tokens and the share of turns with several calls in
\SthreeEOneSinsideR{} of \SthreeEOneSinsideRof{} comparisons (\SthreeEOneCases{} cases, \SthreeEOneMeasures{}
measures; Tables~\ref{tab:s3e1} and~\ref{tab:s3e1tokens}). This holds for S against R and for these measures only.
There was little to organise: without the strategy, \SthreeEOneSeveralRi, \SthreeEOneSeveralRii{} and
\SthreeEOneSeveralRiv{} of the turns that made a call made several (editorial plan, content gaps, billing support;
\SthreeEOneSeveralSi, \SthreeEOneSeveralSii{} and \SthreeEOneSeveralSiv{} with it).

\paragraph{Combined with trimming, the results were mixed, and the hypotheses on the strategy failed as frozen.}
P4-H1 failed in \SthreeHOneFails{} comparison of \SthreeHOneOf: in the editorial plan ST passed the gates in
\SthreeHOneGatesST{} runs against \SthreeHOneGatesT{} for T (the run that failed is the one that asked a question).
P4-H2 failed: on billing support the median turns of S were \SthreeHTwoS, not below R's minimum of \SthreeHTwoR.
P4-H4 failed on billing support, where the median billed input of ST (\SthreeEOneBilledSTiv{} tokens) was above T's
maximum (\SthreeEOneBilledTmaxiv); in the editorial plan it was below T's minimum (\SthreeEOneBilledSTi{} against
\SthreeEOneBilledTmini). With the strategy, the values the prompts prescribe (result limits, search depth, keywords
per call) were \SthreeEOnePrescribedKept{} in all \SthreeEOnePrescribedRuns{} runs of the two cases whose prompts
prescribe them. Runs that asked a question instead of doing the task: \SthreeEOneAsked{} of the
\SthreeEOneStrategyRuns{} with the strategy (in arm ST), \SthreeEOneAskedOther{} of those without it; that run is
kept in every figure. In the billing audit (Section~\ref{sec:s3audit}) the period-end ticket failed in
\AuditEOnePeriod{} runs, all with the strategy text.

\subsection{E3: the strategy text on a second model, without trimming}
\label{sec:s3e3}

E1 left open what the strategy does for a model that issues calls one at a time. E3 asked it of
\texttt{\SthreeModelB}, with arms R and S and no trimming, on billing support and CI triage, whose simulated
servers were live, and on the editorial plan and the content-gap analysis, whose search-data provider was replaced
by recorded outputs: a call got the recorded output of the same call when there was one and otherwise the output of
the recorded call of the same tool with the most similar arguments. The provider was called in \SthreeProviderCallsEThree{} run.
On those two cases only the organisation of calls is evaluated, not the answers. Runs per arm go from
\SthreeEThreeRunsPerArmMin{} to \SthreeEThreeRunsPerArmMax{} (Table~\ref{tab:s3e3}).

\begin{table}[t]
\centering\small
\resizebox{\textwidth}{!}{\begin{tabular}{llrlllrr}
\toprule
Case & Arm & Runs & Turns & Tool calls & Turns with several calls & Billed input & Cost of a run \\
\midrule
Billing support$^\dagger$ & without & 8 & 3.5 [3--8] & 11.5 [10--14] & 88\% [43--100] & 25{,}316 & -- \\
 & strategy & 8 & 3 [3--3] & 10 [10--12] & 100\% [100--100] & 17{,}196 & $-$20\% \\
\addlinespace
CI triage$^\dagger$ & without & 6 & 17 [5--19] & 16 [13--18] & 0\% [0--100] & 865{,}862 & -- \\
 & strategy & 6 & 6 [5--8] & 16.5 [15--20] & 90\% [57--100] & 391{,}018 & $-$15\% \\
\addlinespace
Content gaps$^\ddagger$ & without & 3 & 6 [6--7] & 15 [14--16] & 100\% [100--100] & 169{,}493 & -- \\
 & strategy & 3 & 5 [4--7] & 14 [13--19] & 100\% [100--100] & 131{,}489 & +1\% \\
\addlinespace
Editorial plan$^\ddagger$ & without & 2 & 14 [13--15] & 23 [21--25] & 47\% [36--58] & 1{,}023{,}210 & -- \\
 & strategy & 2 & 10.5 [10--11] & 18 [17--19] & 53\% [50--56] & 658{,}204 & $-$2\% \\
\bottomrule
\end{tabular}
}
\caption{E3, \texttt{\SthreeModelB}, without and with the strategy text, no trimming. Turns, tool calls and turns with several calls as in
Table~\ref{tab:s3e1}; billed input as in Table~\ref{tab:s3e1tokens}; LLM API cost per run relative to the arm without
the strategy, from billed usage at the model's list prices. Billed tokens are
not comparable with E1's (another tokenizer, price list and caching). $^\dagger$Simulated tools, live.
$^\ddagger$Recorded outputs in place of the provider: calls and turns only.}
\label{tab:s3e3}
\end{table}

\paragraph{The hypotheses on turns and grouped calls failed as stated.} P4-H8 (median turns with the strategy below
the fewest turns of any run without it) held in \SthreeHEightHolds{} cases of \SthreeHEightOf, the two with
recorded outputs; P4-H9 (median share of turns with several calls above the highest share of any run without it) held
in \SthreeHNineHolds{} case: in each, a run without the strategy reached the same share or more. Why they failed
concerns the rule, not the runs, whose results stand as they are. Comparing the median of one arm with the extreme
of the other lets a single run decide: with \SthreeEThreeRunsCv{} runs per arm, as on CI triage, one unusual run without the strategy
is enough to fail the hypothesis, and there that is what happened. A protocol that wants to know whether an
arm is usually shorter should compare what is usual in both. The two simulated cases are therefore read apart.

\paragraph{CI triage, where this model issued one call per turn.} Without the strategy, \SthreeEThreeLongRunsv{}
runs of \SthreeEThreeRunsCv{} made every call in a turn of its own (\SthreeEThreeLongMinv{} to
\SthreeEThreeLongMaxv{} turns); the remaining one grouped its calls and took \SthreeEThreeTurnsRminv. With the strategy all
runs took \SthreeEThreeTurnsSminv{} to \SthreeEThreeTurnsSmaxv{} turns, with a similar number of calls. The median
went from \SthreeEThreeTurnsRv{} turns to \SthreeEThreeTurnsSv, billed input tokens fell by
\SthreeEThreeInputChangev, and the LLM API cost per run by \SthreeEThreeCostv: less than the tokens, because most of the
input re-read at each turn is billed at the cached price. P4-H8 fails on this case because of that one run: the
strategy's median, \SthreeEThreeTurnsSv, is not below \SthreeEThreeTurnsRminv.

\paragraph{Billing support, where it already grouped its calls.} Little was left: runs without the strategy took
\SthreeEThreeTurnsRminiv{} to \SthreeEThreeTurnsRmaxiv{} turns, those with it always \SthreeEThreeTurnsSiv, which
is the fewest this task needs (see the post-hoc analysis below).

\paragraph{The implemented gates and coverage checks detected no loss, and the prompts were not overridden.}
P4-H10 holds: the gates passed in
\SthreeHTenGatesiv{} and \SthreeHTenGatesv{} runs of both arms of billing support and CI triage, every requirement of
the prompts covered, as far as the gates and the coverage checks test them. P4-H11 holds: the prescribed values were respected in \SthreeHElevenPassed{} of
\SthreeHElevenChecks{} checks, in the \SthreeHElevenRuns{} strategy runs with recorded outputs. Those two cases point
the same way as CI triage (median turns from \SthreeEThreeTurnsRi{} to \SthreeEThreeTurnsSi{} in the editorial
plan, from \SthreeEThreeTurnsRii{} to \SthreeEThreeTurnsSii{} in the content-gap analysis), on
\SthreeEThreeRunsRi{} and \SthreeEThreeRunsRii{} runs per arm and with recorded outputs: the model's call matched a
recorded call exactly in \SthreeExactMini{} to \SthreeExactMaxi{} of the calls of a run in the editorial plan and in
\SthreeExactMinii{} to \SthreeExactMaxii{} in the content-gap analysis, and got the nearest recorded output otherwise.
They say how calls were organised, not whether the task was done well.

\paragraph{Post hoc: calls that could have been made earlier.} This analysis is not in the protocol; it was added
after the runs, to tell a sequence the task imposes from one the model chose. A call whose arguments name something
an earlier result returned (a run, a job, a customer, an invoice) can only follow that result; any other could have
been made at once. The longest chain of such calls gives the fewest turns a run needed. In CI triage
\SthreeDependingv{} of the calls depended on an earlier result, but the chain was short (the runs, then their jobs,
then the logs): \SthreeFewestv{} turns were enough in every run. Without the strategy \SthreeEarlierRv{} of the calls
were made later than they could have been, with it \SthreeEarlierSv; in billing support \SthreeEarlierRiv{} and
\SthreeEarlierSiv, with \SthreeFewestiv{} turns needed. The rule reads identifiers only and cannot know that a model
wanted to read one log before asking for the next: it gives a lower bound on the turns, not what a careful agent
should do.

\subsection{The hosts checked}
\label{sec:s3hosts}

After the runs of Study~3, the proxy was put in front of servers in several hosts, in most checks with a log that
records each request of the host's client and none of its content; the repository's page on hosts states what was verified,
when and how, and everything else as not verified. In the OpenAI Playground, Claude Code (on Linux and on Windows)
and VS Code (Copilot Chat in agent mode), the tools passed and the model called \texttt{tracetrim\_done} at the end
of its run; in the Playground it did so guided by the tool's description alone. \texttt{expand} was verified in the
Playground only: in the other hosts nothing had yet been trimmed. MCP Inspector, which has no model, showed that the
skill is served as the MCP Skills extension asks. None of the hosts with a model loaded the skill as a skill: no
client declared the extension. The Playground's and VS Code's clients asked for neither the skill nor resources; in Claude Code the model
read the skill through the host's tool for MCP resources before its first call, as a resource. Claude Desktop and
Cursor were not verified. These checks show that the path works in those hosts, not how often or how well.
